\documentclass[conference]{IEEEtran}
\IEEEoverridecommandlockouts

\usepackage[utf8]{inputenc}
\usepackage[T1]{fontenc}
\usepackage{amsmath,amssymb,amsfonts}
\usepackage{algorithmic}
\usepackage{graphicx}
\usepackage{textcomp}
\usepackage{xcolor}
\usepackage{booktabs}
\usepackage{url}
\usepackage[hidelinks]{hyperref}

\begin{document}

\title{Etornie: Verifiable On-Chain Custody and\\
Zero-Knowledge Compliance for Tokenized\\
Intellectual Property on Solana}

\author{
\IEEEauthorblockN{\.Ilknur Timur-Tuncer and the Etornie Research Team}
\IEEEauthorblockA{\textit{Etornie} \\
Ruessenstrasse 5, Baar, Canton of Zug, 6340 Switzerland \\
\url{https://etornieagent.xyz}}
}

\maketitle

\begin{abstract}
Etornie is an intellectual-property (IP) protocol that binds legally
meaningful IP records to verifiable, privacy-preserving on-chain state on
Solana. Rather than tokenizing IP as tradeable financial instruments,
Etornie issues non-transferable, soul-bound records whose authenticity,
ownership, and lifecycle are cryptographically attestable while the
underlying documents and identities remain off-chain and protected. The
protocol combines four constructions: (i) a canonical-hash attestation layer
that anchors case metadata and its full lifecycle to append-only on-chain
events; (ii) a soul-bound Token-2022 issuance scheme in which transferability
is disabled at the protocol level through frozen-by-default token accounts
under program-controlled authority; (iii) a Groth16 zero-knowledge layer over
the BN254 curve that lets a holder prove file ownership and query compliance
without revealing the file, the query, or the witness; and (iv) a
sponsored-transaction custody model in which a backend operator funds and
co-signs user-authorized transactions under a strict instruction allowlist
that provably prevents operator-key drainage. We formalize the attestation,
tokenization, proof, and custody constructions, analyze their security under
a defined threat model, and describe how the design reconciles blockchain
immutability with the EU General Data Protection Regulation and the Swiss
Federal Act on Data Protection by keeping only non-identifying commitments
on-chain. We report compute-cost measurements for single and batched proof
verification and describe an AI-assisted filing agent that integrates with
live European intellectual-property offices.
\end{abstract}

\begin{IEEEkeywords}
intellectual property, real-world assets, Solana, zero-knowledge proofs,
Groth16, BN254, Token-2022, soul-bound tokens, verifiable custody, sponsored
transactions, data protection.
\end{IEEEkeywords}

\section{Introduction}

Intellectual property is among the most valuable and least liquid asset
classes in the modern economy. Registering, proving, licensing, and
transacting rights across trademarks, patents, industrial designs, and
copyrights remains slow, jurisdiction-fragmented, document-centric, and
opaque. Ownership history is scattered across national registries and private
files; provenance is difficult to establish; and the chain of custody over a
rights record is rarely independently verifiable.

Public blockchains offer append-only, independently auditable state and are a
natural substrate for provenance and custody. Yet naive tokenization of IP as
freely tradeable non-fungible tokens is both legally hazardous and
technically inappropriate. A transferable token that purports to represent a
registered right invites securities-law ambiguity, decouples the on-chain
artifact from the off-chain legal instrument, and exposes confidential filing
data on a permanent public ledger.

Etornie takes a deliberately different position. Its thesis is captured by a
single design goal: the intellectual-property operating system for the
on-chain era. IP records are anchored on-chain as \emph{non-transferable,
verifiable attestations}, not as tradeable instruments. The on-chain layer
stores only cryptographic commitments and lifecycle events; all
personally identifying and commercially sensitive material stays off-chain
under conventional data-protection controls. Holders can prove properties of
their records, such as ownership of a specific file or the compliant
provenance of an AI-assisted action, in zero knowledge, revealing nothing
beyond the truth of the statement.

\subsection{Contributions}

\begin{itemize}
\item A canonical-hash \emph{attestation protocol} (Section~\ref{sec:att})
that binds IP case metadata and its complete lifecycle to on-chain
program-derived accounts and append-only events, with init-once replay
resistance.
\item A \emph{soul-bound tokenization} scheme (Section~\ref{sec:token}) built
on the SPL Token-2022 standard, in which non-transferability is enforced by
the token runtime itself through frozen-by-default accounts whose freeze
authority is a program-derived address.
\item A \emph{zero-knowledge layer} (Section~\ref{sec:zk}) with three Groth16
circuits over BN254, verified on-chain through Solana pairing syscalls,
supporting private proofs of file ownership and query compliance, plus a
batch verifier that amortizes pairing cost across up to four proofs.
\item A \emph{verifiable sponsored-custody model} (Section~\ref{sec:custody})
in which a backend operator pays fees and co-signs, and a formal instruction
allowlist guarantees that co-signing cannot be abused to move operator funds,
closing a hot-wallet drainage vector inherent to naive sponsorship.
\item A \emph{data-protection analysis} (Section~\ref{sec:compliance})
showing how on-chain immutability coexists with the right to erasure by
confining on-chain state to non-identifying commitments.
\end{itemize}

\section{System Model and Trust Assumptions}
\label{sec:model}

\subsection{Actors}

\begin{itemize}
\item \textbf{Client} $U$: the rights holder. Controls a Solana wallet
$w_c$ and signs the transactions that authorize records about their own IP.
\item \textbf{Operator} $\mathcal{O}$: an Etornie backend key that funds
transactions (pays network fees and rent) and co-signs sponsored
transactions. The operator provides liveness and gas abstraction; it is not a
custodian of the client's rights.
\item \textbf{Programs}: three on-chain Solana programs, for attestation,
tokenization, and zero-knowledge verification, whose bytecode is public and
whose upgrade authority is intended to be held by a multi-signature.
\item \textbf{IP offices}: external registries (for example the European
Union Intellectual Property Office) with which the platform interoperates
through authenticated adapters.
\end{itemize}

\subsection{Trust and threat model}

We assume the Solana ledger provides safety and liveness under its consensus
assumptions, and that BN254 pairings and the SHA-256 and Poseidon hash
functions are secure. The client's wallet key is trusted by the client. The
operator is trusted for \emph{liveness} (it may refuse service) but is
\emph{not} trusted for \emph{integrity}: the protocol is designed so that a
compromised or malicious operator cannot forge a record the client did not
authorize, cannot make the client's soul-bound token transferable, and cannot,
through the sponsorship mechanism, be coerced into signing a transaction that
spends operator or client funds outside the defined fee. Adversaries may
submit arbitrary transactions to the programs and arbitrary payloads to the
backend. We do not defend against a fully compromised client wallet, nor
against off-chain legal disputes over the underlying rights, which remain the
province of the competent registries and courts.

\section{On-Chain Attestation Protocol}
\label{sec:att}

The attestation program anchors an IP case and its evolution to on-chain
state. A case is identified by a 128-bit identifier $\mathrm{cid}$.

\subsection{Metadata commitment}

Let $m$ denote the structured metadata of a case, including its identifier,
type, jurisdiction, classification, filing and deadline dates, and the client
wallet. Etornie serializes $m$ into a canonical form
$\mathrm{canon}(m)$ using lexicographically sorted keys and compact
separators, then commits to it with SHA-256:
\begin{equation}
h_m \;=\; \mathrm{SHA256}\big(\mathrm{canon}(m)\big) \in \{0,1\}^{256}.
\end{equation}
Canonicalization ensures the commitment is independent of serialization
order, so a verifier who reconstructs $m$ can recompute and check $h_m$
deterministically. Only $h_m$, never $m$, reaches the chain.

\subsection{Record and derivation}

The attestation record is a program-derived account (PDA) keyed by the case
identifier:
\begin{equation}
\mathrm{PDA}_{\mathrm{att}} \;=\; \mathcal{D}\big(\texttt{"case"} \,\|\,
\mathrm{cid};\ \mathsf{prog}_{\mathrm{att}}\big),
\end{equation}
where $\mathcal{D}$ is Solana's PDA derivation. The stored record is the tuple
\begin{equation}
R_{\mathrm{att}} = \big(\mathrm{cid},\, h_m,\, w_{\mathrm{cr}},\, w_c,\,
w_{\mathcal{O}},\, t\big),
\end{equation}
binding the case to its metadata commitment $h_m$, the creator signer
$w_{\mathrm{cr}}$, the client wallet $w_c$, the operator $w_{\mathcal{O}}$, and
the ledger timestamp $t$. The account is created with an init-once constraint,
so a given $\mathrm{cid}$ can be attested exactly once; a replayed creation
fails at the runtime level. This yields the uniqueness invariant
\begin{equation}
\forall\, \mathrm{cid}:\ \big|\{R_{\mathrm{att}} : R_{\mathrm{att}}.\mathrm{cid}
= \mathrm{cid}\}\big| \le 1.
\end{equation}

\subsection{Lifecycle events}

State transitions (for example filing, examination, grant, or closure) are
recorded by an update instruction that overwrites $h_m$ with the fresh
commitment and emits a typed event
\begin{equation}
\mathsf{Ev} = \big(\mathrm{cid},\, h_m^{\mathrm{old}},\, h_m^{\mathrm{new}},\,
\eta,\, w_{\mathrm{actor}},\, t\big),
\end{equation}
where $\eta$ is a lifecycle discriminator. Because every transition is an
on-chain event, the complete, tamper-evident history of a case is
reconstructable by replaying the ordered event log
\begin{equation}
\mathcal{T} = \big\langle \mathsf{Ev}_0, \mathsf{Ev}_1, \ldots,
\mathsf{Ev}_k \big\rangle,
\end{equation}
without exposing any case content beyond successive commitments.

\section{Soul-Bound Tokenization}
\label{sec:token}

Once attested, a case may be tokenized. Etornie issues a token that is
\emph{soul-bound}: bound to the client wallet at issuance and
non-transferable thereafter. Transferability is not merely discouraged by
convention; it is disabled by the token runtime.

\subsection{Construction}

The token is an SPL Token-2022 mint with zero decimals and unit supply. Its
mint and freeze authorities are a program-derived authority $\pi$, and the
mint is configured with the \emph{default account state} extension set to
\textsf{Frozen}. Consequently, the client's associated token account $a$ is
created in the frozen state. Minting proceeds as an atomic
thaw-mint-refreeze under the program authority:
\begin{equation}
\pi:\quad \mathsf{Thaw}(a)\ \to\ \mathsf{MintTo}(a, 1)\ \to\
\mathsf{Freeze}(a).
\end{equation}

\subsection{Non-transferability invariant}

After issuance, $a$ is frozen and only $\pi$ can thaw it. The Token-2022
runtime rejects any transfer, delegation, or burn on a frozen account.
Therefore, for any transaction $\tau$ authored by the client,
\begin{equation}
\mathrm{state}(a) = \textsf{Frozen} \ \wedge\ \mathsf{FreezeAuth}(a) = \pi
\;\Rightarrow\; \mathsf{Transfer}(a) \notin \tau_{\mathrm{valid}} .
\end{equation}
Burning is likewise gated: only an operator-authorized program instruction
can thaw, burn one unit, and mark the record burned, with an
already-burned guard preventing replay. The token thus functions as a
verifiable, revocable certificate of record rather than a tradeable asset,
which is the property that keeps it outside the ambit of transferable
financial instruments.

\section{Zero-Knowledge Proof Layer}
\label{sec:zk}

Etornie proves properties of off-chain artifacts without revealing them,
using Groth16 arguments \cite{groth16} over the BN254 pairing-friendly curve
\cite{bn254}, verified on-chain through Solana's \texttt{alt\_bn128}
syscalls by the Etornie Mosaic verifier.

\subsection{Groth16 verification}

For a circuit with public input vector $x = (x_1, \ldots, x_\ell)$ and proof
$\pi = (A, B, C)$ with $A, C \in \mathbb{G}_1$ and $B \in \mathbb{G}_2$,
verification checks the pairing equation
\begin{equation}
e(A, B) \;=\; e(\alpha, \beta)\cdot e\!\left(\sum_{i=0}^{\ell} x_i\,
\mathrm{IC}_i,\ \gamma\right)\cdot e(C, \delta),
\label{eq:groth}
\end{equation}
with $x_0 = 1$ and verifying-key elements $\alpha, \beta, \gamma, \delta,
\{\mathrm{IC}_i\}$. Writing the public-input linear combination as
\begin{equation}
L \;=\; \mathrm{IC}_0 + \sum_{i=1}^{\ell} x_i\, \mathrm{IC}_i \in \mathbb{G}_1,
\end{equation}
\eqref{eq:groth} is equivalent to the single multi-pairing identity
\begin{equation}
e(-A, B)\cdot e(\alpha, \beta)\cdot e(L, \gamma)\cdot e(C, \delta)
\;=\; 1_{\mathbb{G}_T},
\label{eq:multipair}
\end{equation}
which the verifier evaluates in one batched \texttt{alt\_bn128} pairing
syscall. A proof is $256$ bytes ($64\,\|\,128\,\|\,64$), and the verifying key
is serialized canonically as
$\alpha\,\|\,\beta\,\|\,\gamma\,\|\,\delta\,\|\,\{\mathrm{IC}_i\}$ with
$\ell + 1$ input-commitment points.

\subsection{Statements}

Two application circuits share a commitment structure. A 256-bit content hash
$h$ (a SHA-256 digest of a file or of an AI query) is split into two 128-bit
halves $h_{\mathrm{hi}}, h_{\mathrm{lo}}$, each embedded in a BN254 field
element, and bound to a Poseidon \cite{poseidon} commitment of a secret
witness $s$:
\begin{align}
\mathcal{R}_{\mathrm{own}} &= \big\{ (h_{\mathrm{hi}}, h_{\mathrm{lo}},
\mathrm{cm});\, s :\ \mathrm{cm} = \mathsf{Poseidon}(s, h_{\mathrm{hi}},
h_{\mathrm{lo}}) \big\}, \\
\mathcal{R}_{\mathrm{cmp}} &= \big\{ (q_{\mathrm{hi}}, q_{\mathrm{lo}},
\mathrm{cm});\, s :\ \mathrm{cm} = \mathsf{Poseidon}(s, q_{\mathrm{hi}},
q_{\mathrm{lo}}) \big\}.
\end{align}
Relation $\mathcal{R}_{\mathrm{own}}$ is a proof of \emph{file ownership}: the
prover demonstrates knowledge of a secret $s$ that commits to a specific file
hash, without revealing $s$ or the file. Relation $\mathcal{R}_{\mathrm{cmp}}$
is a proof of \emph{query compliance} for AI-assisted actions, where the
secret is derived from a wallet signature off-chain. Zero knowledge keeps the
file, the query, and the witness private; soundness prevents a party from
producing a valid commitment for content whose witness it does not know. On
success the verifier initializes a record PDA keyed by
$(\text{domain}, w_c, h)$, giving per-holder, per-artifact replay resistance.

\subsection{Payment binding}

For compliance proofs that gate a paid action, the on-chain payment is bound
to the specific proof. The paying transfer carries a memo
\begin{equation}
\mu \;=\; \mathsf{Base58}\big(\mathrm{SHA256}(q \,\|\, \mathrm{cm})\big),
\end{equation}
where $q$ is the query-hash encoding and $\mathrm{cm}$ the commitment. A
payment is accepted only if its memo matches the proof under verification,
preventing a valid payment from being reused for a different query or a valid
proof from being settled by an unrelated payment.

\subsection{Batch verification}

For throughput, up to $n \in [2, 4]$ proofs are verified together by a batch
verifier that aggregates their pairing checks into a single evaluation, at the
cost of one field element of randomness per proof. Each batch entry is
$288$ bytes ($256$-byte proof and a $32$-byte journal digest), so four entries
fit within Solana's transaction size limit. The batch is keyed by
\begin{equation}
d = \mathrm{SHA256}\Big(\big\|_{j=1}^{n}\, \mathrm{SHA256}\big(\big\|_i\,
x_i^{(j)}\big)\Big),
\end{equation}
recomputed on-chain to prevent seed spoofing. Batching reduces the marginal
verification cost from approximately $83{,}500$ to approximately $52{,}000$
compute units per proof, a reduction of about $38\%$
(Table~\ref{tab:cost}).

\section{Verifiable Sponsored Custody and Fee Enforcement}
\label{sec:custody}

A central usability requirement is that clients need not hold or spend SOL to
create records. Etornie therefore sponsors transactions: the operator is the
fee payer, and the client contributes only the signature that authorizes an
action about their own IP.

\subsection{Co-signing}

A Solana transaction has an ordered list of required signers; the first is the
fee payer. In Etornie's flow the client's application assembles a versioned
transaction, the client wallet signs its slot, and the backend splices the
operator signature into the fee-payer slot before submission. This gives the
client gas abstraction while preserving the property that the client's
authorization is a first-class signature over the exact message.

\subsection{The drainage hazard and the allowlist}

Naive sponsorship is dangerous: an operator that blindly co-signs whatever a
client submits can be induced to authorize instructions that spend operator
funds. Concretely, an adversarial client could submit a transaction containing
both the expected instruction and a system transfer
$\mathcal{O} \to \mathrm{adv}$; because the operator is a signer, the transfer
would execute and drain the operator's hot wallet.

Etornie closes this vector by verifying the full instruction set before
co-signing. Let $\tau$ be the submitted transaction, $\mathcal{P}$ the set of
program identifiers expected for the flow, $\mathsf{Sys}$ the system program,
$\Phi$ the treasury, and $\phi$ the required fee. Define the safety predicate
\begin{equation}
\mathsf{Safe}(\tau) =
\underbrace{\neg\,\mathsf{LUT}(\tau)}_{\text{no lookup tables}}
\ \wedge\
\underbrace{\bigwedge_{I \in \tau}\ \mathrm{prog}(I) \in \mathcal{P} \cup
\{\mathsf{Sys}\}}_{\text{program allowlist}}
\ \wedge\
\mathsf{Fee}(\tau),
\label{eq:safe}
\end{equation}
where the fee clause requires that every system instruction be exactly a
transfer to the treasury sourced from a non-operator signer, and that the
total so transferred meets the fee:
\begin{equation}
\mathsf{Fee}(\tau) = \Big(\!\!\sum_{\substack{I \in \tau,\ I = \mathsf{tr}(u
\to \Phi)\\ u \in \mathsf{Signers}(\tau)\setminus\{\mathcal{O}\}}}\!\!
\mathrm{lamports}(I)\ \ge\ \phi\Big).
\end{equation}
The operator signs $\tau$ if and only if $\mathsf{Safe}(\tau)$ holds.
Rejecting address lookup tables prevents accounts, including the treasury or
the operator, from being resolved dynamically to evade the static check;
restricting programs to $\mathcal{P} \cup \{\mathsf{Sys}\}$ blocks invocation
of unexpected programs under the operator signature; and constraining every
system instruction to be the user-to-treasury fee blocks the drainage
transfer. Thus a compromised or adversarial client cannot spend operator
funds, and the fee is non-bypassable: a client that omits the transfer fails
$\mathsf{Fee}$, and a client that adds a hostile instruction fails the
allowlist.

\subsection{Custody guarantee}

Because the client's own signature authorizes each record, and the operator's
signature is constrained by \eqref{eq:safe} to fee payment plus the expected
program instruction, the operator can neither forge a record nor divert funds.
The operator is reduced to a liveness and gas provider, which is the intended
minimal-trust custody posture.

\section{Economic Model}
\label{sec:econ}

Etornie's on-chain economics are denominated in SOL and collected into a
treasury $\Phi$. Table~\ref{tab:fees} lists the protocol fees. A registration
fee is charged at attestation, a mint fee at tokenization, and a micro-fee
per AI query settled through the compliance-proof and payment-binding path of
Section~\ref{sec:zk}. All fees are enforced by the custody predicate of
Section~\ref{sec:custody} and are inert until a treasury is configured, so the
mechanism is safe by default in development.

\begin{table}[t]
\caption{Protocol fees}
\label{tab:fees}
\centering
\begin{tabular}{@{}lrl@{}}
\toprule
Action & Fee (SOL) & Enforcement \\
\midrule
Registration (attestation) & $0.01$ & bundled, allowlisted \\
Tokenization (mint) & $0.05$ & bundled, allowlisted \\
AI query (compliance) & $0.0001$ & payment-bound proof \\
\bottomrule
\end{tabular}
\end{table}

At launch the protocol has no native token; value accrues to the treasury in
SOL per registered and tokenized asset. A native protocol token is planned to
support decentralized governance and incentive alignment among filers,
attestation validators, and the treasury. Its design, distribution, and any
utility are outside the scope of this paper and will be specified separately.
Nothing in this document constitutes an offer of, or solicitation for, any
security or financial instrument.

\section{AI-Assisted Filing}
\label{sec:agent}

Etornie exposes an agentic interface, EtornieGPT, that assists rights holders
through search, drafting, fee estimation, and submission. The agent operates
over tool interfaces to external IP offices through authenticated adapters. At
present the European Union Intellectual Property Office is integrated as a
live OAuth2 REST client for trademark search and electronic filing, and the
United Kingdom Intellectual Property Office is integrated through a supervised
browser-automation adapter that drives the official electronic form up to the
payment step. Additional offices, including the World Intellectual Property
Organization Global Brand Database, are planned through the same adapter
abstraction. Every AI query that gates a paid action is settled through the
compliance-proof mechanism of Section~\ref{sec:zk}, so the agent's paid
actions are both metered and privately attestable.

\section{Compliance and Data Protection}
\label{sec:compliance}

Etornie processes personal and commercially sensitive data and operates from
Switzerland, and is therefore designed against both the EU General Data
Protection Regulation \cite{gdpr} and the revised Swiss Federal Act on Data
Protection \cite{fadp}. The central tension between these regimes and a public
ledger is the right to erasure versus on-chain immutability.

Etornie resolves the tension by construction. The on-chain layer stores only
\emph{non-identifying commitments}: SHA-256 metadata hashes, Poseidon
commitments, content-hash halves, and wallet public keys, none of which
reveals the underlying documents, identities, or filing content. All personal
and confidential material remains off-chain in conventional storage subject to
access control, retention limits, and deletion. Because the on-chain
commitments are one-way and content-hiding, erasing the off-chain data renders
the on-chain hashes non-identifying and inert, so a data subject's erasure
request can be honored off-chain without any conflict with ledger
immutability. The soul-bound design further limits exposure: records are not
tradeable, so IP data is never propagated through a secondary market.

\section{Security Analysis}
\label{sec:security}

We summarize the guarantees against the threat model of
Section~\ref{sec:model}.

\emph{Record integrity.} A record requires the client's signature; the
operator cannot forge one. Attestation is init-once, so records cannot be
overwritten or replayed for a given case identifier.

\emph{Non-transferability.} Soul-binding is enforced by the token runtime
through a frozen account under a program freeze authority, not by off-chain
policy; a malicious operator cannot make a token transferable without the
program authority.

\emph{Proof soundness and privacy.} Groth16 soundness prevents forging proofs
for statements without a valid witness; zero knowledge hides files, queries,
and witnesses. Per-holder, per-artifact record PDAs and on-chain recomputation
of batch digests prevent proof replay and seed spoofing.

\emph{Custody safety.} The allowlist predicate \eqref{eq:safe} guarantees the
operator co-signs only fee payment plus the expected program instruction,
eliminating the hot-wallet drainage vector and making fees non-bypassable.

\emph{Residual assumptions.} We rely on the security of BN254 pairings,
SHA-256, and Poseidon; on the correctness of the deployed program bytecode,
for which independent audit is planned before mainnet; and on the client
protecting its own wallet key.

\section{Related Work}
\label{sec:related}

On-chain attestation frameworks provide general schemas for making signed
claims about arbitrary subjects; Etornie specializes attestation to IP
lifecycles with canonical-hash commitments and program-enforced uniqueness.
Non-transferable tokens, or soul-bound tokens \cite{sbt}, motivate
identity- and credential-bound assets; Etornie realizes the property at the
token-runtime level rather than through social convention, using Token-2022
account freezing. Succinct non-interactive arguments, and Groth16
\cite{groth16} in particular, underpin our proof layer, with Poseidon
\cite{poseidon} providing an arithmetization-friendly commitment and BN254
\cite{bn254} enabling native on-chain pairing verification. Relative to
real-world-asset tokenization that emphasizes tradeable representations,
Etornie deliberately forgoes transferability in favor of verifiable custody
and privacy, which we argue is the appropriate posture for regulated
intellectual-property records.

\section{Conclusion and Future Work}
\label{sec:conclusion}

Etornie shows that intellectual property can be given verifiable, private,
and legally coherent on-chain representation without turning rights into
tradeable tokens. Canonical-hash attestation anchors provenance and
lifecycle; soul-bound Token-2022 issuance binds records to their holders;
a Groth16 layer proves ownership and compliance in zero knowledge; and a
formally constrained sponsorship model provides gas abstraction without
custodial risk, while keeping only non-identifying commitments on-chain for
data-protection compliance. Future work includes independent security audit,
a multi-signature upgrade authority, expansion of the IP-office adapter set,
validated priority-fee support for mainnet transaction landing, and a
separately specified governance token.

\section*{Disclaimer}
This document is a technical description for informational purposes only. It
does not constitute legal, financial, or investment advice, nor an offer or
solicitation of any security or financial instrument. Nothing herein creates a
representation as to the legal status of any right in any jurisdiction, which
remains subject to the competent registries and courts.

\begin{thebibliography}{99}

\bibitem{groth16}
J. Groth, ``On the size of pairing-based non-interactive arguments,'' in
\emph{Advances in Cryptology (EUROCRYPT)}, 2016, pp. 305--326.

\bibitem{bn254}
P. S. L. M. Barreto and M. Naehrig, ``Pairing-friendly elliptic curves of
prime order,'' in \emph{Selected Areas in Cryptography (SAC)}, 2005, pp.
319--331.

\bibitem{poseidon}
L. Grassi, D. Khovratovich, C. Rechberger, A. Roy, and M. Schofnegger,
``Poseidon: A new hash function for zero-knowledge proof systems,'' in
\emph{USENIX Security Symposium}, 2021, pp. 519--535.

\bibitem{sbt}
E. G. Weyl, P. Ohlhaver, and V. Buterin, ``Decentralized society: Finding
Web3's soul,'' SSRN Working Paper 4105763, 2022.

\bibitem{solana}
A. Yakovenko, ``Solana: A new architecture for a high performance
blockchain,'' Whitepaper, 2018.

\bibitem{token2022}
Solana Program Library, ``Token-2022 program and extensions,''
\url{https://spl.solana.com/token-2022}, accessed 2026.

\bibitem{gdpr}
European Parliament and Council, ``Regulation (EU) 2016/679 (General Data
Protection Regulation),'' \emph{Official Journal of the European Union}, 2016.

\bibitem{fadp}
Swiss Confederation, ``Federal Act on Data Protection (revised FADP),'' in
force 1 September 2023.

\end{thebibliography}

\end{document}
